\documentclass[aspectratio=169]{beamer}
\usepackage[english]{babel}
\usepackage[utf8]{inputenc}
\usepackage[T1]{fontenc}
\usepackage{lmodern}
\usepackage{listings}
\usepackage{xcolor}
\usepackage{graphicx}
\usepackage{textcomp}
\DeclareUnicodeCharacter{2014}{---}
\DeclareUnicodeCharacter{2013}{--}
\DeclareUnicodeCharacter{2212}{-}
\DeclareUnicodeCharacter{2192}{\ensuremath{\rightarrow}}
\DeclareUnicodeCharacter{00B1}{\ensuremath{\pm}}
\DeclareUnicodeCharacter{00D7}{\ensuremath{\times}}
\DeclareUnicodeCharacter{00B7}{\ensuremath{\cdot}}
\DeclareUnicodeCharacter{20AC}{EUR}
\DeclareUnicodeCharacter{2070}{\textsuperscript{0}}
\DeclareUnicodeCharacter{00B9}{\textsuperscript{1}}
\DeclareUnicodeCharacter{00B2}{\textsuperscript{2}}
\DeclareUnicodeCharacter{00B3}{\textsuperscript{3}}
\DeclareUnicodeCharacter{2074}{\textsuperscript{4}}
\DeclareUnicodeCharacter{2075}{\textsuperscript{5}}
\DeclareUnicodeCharacter{2076}{\textsuperscript{6}}
\DeclareUnicodeCharacter{2077}{\textsuperscript{7}}
\DeclareUnicodeCharacter{2078}{\textsuperscript{8}}
\DeclareUnicodeCharacter{2079}{\textsuperscript{9}}
\DeclareUnicodeCharacter{2248}{\ensuremath{\approx}}
\DeclareUnicodeCharacter{2264}{\ensuremath{\le}}

% Colors for code
\definecolor{codebg}{RGB}{245,245,245}
\definecolor{codecomment}{RGB}{0,128,0}
\definecolor{codekeyword}{RGB}{0,0,255}
\definecolor{codestring}{RGB}{163,21,21}
\definecolor{bandcolor}{RGB}{200,50,50}

\lstdefinestyle{pythonstyle}{
    language=Python,
    backgroundcolor=\color{codebg},
    basicstyle=\ttfamily\scriptsize,
    keywordstyle=\color{codekeyword},
    commentstyle=\color{codecomment},
    stringstyle=\color{codestring},
    showstringspaces=false,
    breaklines=true,
    frame=single,
    rulecolor=\color{black},
    escapeinside={(*@}{@*)},
    moredelim=**[l][\color{bandcolor}\bfseries]{\#\ ---\ NEW},
    moredelim=**[l][\color{bandcolor}\bfseries]{\#\ ---\ CHANGED},
    literate=
      {á}{{\'a}}1 {é}{{\'e}}1 {í}{{\'i}}1 {ó}{{\'o}}1 {ú}{{\'u}}1
      {Á}{{\'A}}1 {É}{{\'E}}1 {Í}{{\'I}}1 {Ó}{{\'O}}1 {Ú}{{\'U}}1
      {ñ}{{\~n}}1 {Ñ}{{\~N}}1 {ü}{{\"u}}1 {Ü}{{\"U}}1
      {¿}{{?`}}1 {¡}{{!`}}1
      {—}{{---}}1 {–}{{--}}1 {−}{{-}}1
      {±}{{\ensuremath{\pm}}}1 {×}{{\ensuremath{\times}}}1
      {→}{{\ensuremath{\rightarrow}}}1
      {°}{{\textdegree}}1
      {·}{{\ensuremath{\cdot}}}1
      {€}{{\texteuro}}1
      {≈}{{\ensuremath{\approx}}}1
      {≤}{{\ensuremath{\le}}}1
      {⁰}{{\textsuperscript{0}}}1 {¹}{{\textsuperscript{1}}}1
      {²}{{\textsuperscript{2}}}1 {³}{{\textsuperscript{3}}}1
      {⁴}{{\textsuperscript{4}}}1 {⁵}{{\textsuperscript{5}}}1
      {⁶}{{\textsuperscript{6}}}1 {⁷}{{\textsuperscript{7}}}1
      {⁸}{{\textsuperscript{8}}}1 {⁹}{{\textsuperscript{9}}}1
}

\lstset{style=pythonstyle}

% Title slide macro: \lessontitle{Lesson Number}{Title}
\newcommand{\lessontitle}[2]{
    \title{Lesson #1: #2}
    \author{}
    \institute{Tec de Monterrey}
    \begin{frame}[plain]
        \titlepage
    \end{frame}
}

% Figure frame macro: \figureframe{Title}{Image path}
\newcommand{\figureframe}[2]{
    \begin{frame}{#1}
        \begin{center}
            \includegraphics[height=0.8\textheight,width=\textwidth,keepaspectratio]{#2}
        \end{center}
    \end{frame}
}


% Listings pulled from src/ with \lstinputlisting, so the slide is the file.
\lstset{
    basicstyle=\ttfamily\fontsize{8pt}{9.6pt}\selectfont,
    breaklines=true,
    breakatwhitespace=true,
    columns=fullflexible,
    keepspaces=true,
    xleftmargin=0.3em,
    framesep=2pt,
    aboveskip=0pt,
    belowskip=0pt
}

\newcommand{\resultframe}[3]{%
    \begin{frame}{#1}
        \begin{center}
            \includegraphics[height=0.62\textheight,width=\textwidth,keepaspectratio]{#2}
        \end{center}
        \vspace{0.25em}
        {\small #3\par}
    \end{frame}%
}

\newcommand{\claimframe}[2]{%
    \begin{frame}{#1}
        {\small #2\par}
    \end{frame}%
}

\date{}

\begin{document}

\lessontitle{08}{BlackBox}

\section{This lesson}

\begin{frame}{This lesson}
Lesson 07 tuned knobs on a landscape we could still see. This function arrives from somewhere else: five arguments, three kinds, no honest plot.

\vspace{0.6em}
\begin{itemize}
    \item A black box states legal inputs and says nothing about the output.
    \item Mixed real, integer and categorical genes need explicit legality.
    \item Twelve searches agreeing on \(x\) while fitness differs by thousands is not a maximum.
\end{itemize}
\end{frame}

% ================= black_box_01_no_picture =================
\begin{frame}{A mixed domain, and a dangerous objective}
\[
\mathcal X=[a_l,a_u]\times[b_l,b_u]\times[x_l,x_u]
\times\mathcal N\times\mathcal F.
\]
Constructor repair projects proposals back to this Cartesian product and changes
their distribution. Here
\(D=(n+1)^2(1+a+b)(120-x^2)r_\phi(x;b,n)+\tfrac12\) changes sign.
Bisection locates \(D(x^*)=0\). That point is a pole, so the maximization problem
is unbounded or ill-posed on a domain containing \(x^*\), not ``solved''.
\end{frame}

\section{The function you cannot look at}

\begin{frame}{The function you cannot look at}
Every lesson so far has optimised f(x) = sin(x) - 0.2*|x| or a close relative:
one real gene, one curve, one picture on the slide, and - crucially - a
brute-force optimum on a grid to check the run against. Lesson 02 step 5 said so
in as many words and warned that the honest answer to "did this run succeed?"
gets hard once x stops being one number.
\end{frame}

\resultframe{The function you cannot look at}{../figures/black_box_01_no_picture.png}{%
    Five live slices agree on \texttt{x = -10.950}, but their maxima differ by \texttt{9.23×}; the picture has not resolved the value}

% ================= black_box_02_the_grid =================
\section{The grid, priced and then refined}

\begin{frame}{The grid, priced and then refined}
Every earlier lesson checked its runs against a brute-force optimum on a grid.
That check is what made "did this run succeed?" answerable. This step tries to
buy it here, and fails twice.
\end{frame}

\begin{frame}[fragile]{(1) evaluation\_cost()}
\lstinputlisting[basicstyle=\ttfamily\fontsize{7pt}{8.5pt}\selectfont,firstline=67,lastline=87]{../src/black_box_02_the_grid.py}
\end{frame}

\begin{frame}[fragile]{(2) BOOK\_GRID}
\lstinputlisting[basicstyle=\ttfamily\fontsize{8pt}{9.6pt}\selectfont,firstline=89,lastline=101]{../src/black_box_02_the_grid.py}
\end{frame}

\begin{frame}[fragile]{(3) grid\_maximum()}
\lstinputlisting[basicstyle=\ttfamily\fontsize{7pt}{8.5pt}\selectfont,firstline=120,lastline=143]{../src/black_box_02_the_grid.py}
\end{frame}

\begin{frame}[fragile]{(4) REFINEMENTS}
\lstinputlisting[basicstyle=\ttfamily\fontsize{8pt}{9.6pt}\selectfont,firstline=145,lastline=154]{../src/black_box_02_the_grid.py}
\end{frame}

\resultframe{The grid, priced and then refined}{../figures/black_box_02_the_grid.png}{%
    The book grid costs 105,525,000 calls; refinement makes the reported maximum grow 61,099-fold instead of converge}

% ================= black_box_03_the_chromosome =================
\section{One chromosome, three kinds of gene}

\begin{frame}{One chromosome, three kinds of gene}
Until now a chromosome was a real number, or a vector of real numbers, and
"legal" meant "inside [-10, 10]". Here one individual carries three different
kinds of gene at once:
\end{frame}

\begin{frame}[fragile]{(1) clamp()}
\lstinputlisting[basicstyle=\ttfamily\fontsize{8pt}{9.6pt}\selectfont,firstline=80,lastline=89]{../src/black_box_03_the_chromosome.py}
\end{frame}

\begin{frame}[fragile]{(2) closest()}
\lstinputlisting[basicstyle=\ttfamily\fontsize{8pt}{9.6pt}\selectfont,firstline=91,lastline=102]{../src/black_box_03_the_chromosome.py}
\end{frame}

\begin{frame}[fragile]{(3) Individual}
\lstinputlisting[basicstyle=\ttfamily\fontsize{6pt}{7.2pt}\selectfont,firstline=104,lastline=134]{../src/black_box_03_the_chromosome.py}
\end{frame}

\begin{frame}[fragile]{(4) create\_random()}
\lstinputlisting[basicstyle=\ttfamily\fontsize{8pt}{9.6pt}\selectfont,firstline=136,lastline=153]{../src/black_box_03_the_chromosome.py}
\end{frame}

\begin{frame}[fragile]{(5) repair\_report()}
\lstinputlisting[basicstyle=\ttfamily\fontsize{8pt}{9.6pt}\selectfont,firstline=155,lastline=171]{../src/black_box_03_the_chromosome.py}
\end{frame}

\resultframe{One chromosome, three kinds of gene}{../figures/black_box_03_the_chromosome.png}{%
    The constructor repairs all five illegal examples; the initial population samples only 10\% of the declared \texttt{x} range}

% ================= black_box_04_the_operators =================
\section{The operators, and what silent repair costs}

\begin{frame}{The operators, and what silent repair costs}
Step 3 put all legality in the constructor, so the operators can be written
as if every gene were a real number. This step writes them - the book's
design, and a deliberate one: every gene gets its own coin and its own step
size, because the genes are different kinds of object. Blend alpha 0.3 for
a, 0.5 for b, 1.0 for x; gaussian sigma 0.2 for a and b, 1.0 for x and n;
the label gene mutates by re-choosing, the only mutation a set admits.
\end{frame}

\begin{frame}[fragile]{(1) crossover()}
\lstinputlisting[basicstyle=\ttfamily\fontsize{7pt}{8.5pt}\selectfont,firstline=126,lastline=151]{../src/black_box_04_the_operators.py}
\end{frame}

\begin{frame}[fragile]{(2) mutate()}
\lstinputlisting[basicstyle=\ttfamily\fontsize{7pt}{8.5pt}\selectfont,firstline=153,lastline=174]{../src/black_box_04_the_operators.py}
\end{frame}

\begin{frame}[fragile]{(3) the repair census}
\lstinputlisting[basicstyle=\ttfamily\fontsize{6pt}{7.2pt}\selectfont,firstline=192,lastline=223]{../src/black_box_04_the_operators.py}
\end{frame}

\begin{frame}[fragile]{(4) the drift demo}
\lstinputlisting[basicstyle=\ttfamily\fontsize{7pt}{8.5pt}\selectfont,firstline=225,lastline=248]{../src/black_box_04_the_operators.py}
\end{frame}

\resultframe{The operators, and what silent repair costs}{../figures/black_box_04_the_operators.png}{%
    41.5\% of forced integer mutations snap back unchanged; coarse lattices can silence mutation completely}

% ================= black_box_05_the_search =================
\section{The search, and an answer that will not sit still}



\begin{frame}[fragile]{(1) select\_rank\_with\_elite()}
\lstinputlisting[basicstyle=\ttfamily\fontsize{7pt}{8.5pt}\selectfont,firstline=166,lastline=191]{../src/black_box_05_the_search.py}
\end{frame}

\begin{frame}[fragile]{(2) the generational loop}
\lstinputlisting[basicstyle=\ttfamily\fontsize{7pt}{8.5pt}\selectfont,firstline=194,lastline=219]{../src/black_box_05_the_search.py}
\end{frame}

\begin{frame}[fragile]{(3) the diagnostics}
\lstinputlisting[basicstyle=\ttfamily\fontsize{7pt}{8.5pt}\selectfont,firstline=229,lastline=254]{../src/black_box_05_the_search.py}
\end{frame}

\resultframe{The search, and an answer that will not sit still}{../figures/black_box_05_the_search.png}{%
    The best fitness reaches \texttt{2.278e-09}, but the curve is still rising and 97.5\% of the final population has left the initial \texttt{x} interval}

% ================= black_box_06_the_verdict =================
\section{The verdict, twelve times over}

\begin{frame}{The verdict, twelve times over}
Step 5's run looked like a success, and one run is an anecdote. This step
repeats the whole search twelve times, from twelve seeds, and lays the
twelve champions side by side. Nothing about the algorithm changes: same
chromosome, same operators, same knobs, same 100 generations. Only the seed
moves.
\end{frame}

\begin{frame}[fragile]{(1) run\_search()}
\lstinputlisting[basicstyle=\ttfamily\fontsize{7pt}{8.5pt}\selectfont,firstline=180,lastline=203]{../src/black_box_06_the_verdict.py}
\end{frame}

\begin{frame}[fragile]{(2) the Monte Carlo}
\lstinputlisting[basicstyle=\ttfamily\fontsize{8pt}{9.6pt}\selectfont,firstline=208,lastline=225]{../src/black_box_06_the_verdict.py}
\end{frame}

\begin{frame}[fragile]{(3) the two readings}
\lstinputlisting[basicstyle=\ttfamily\fontsize{8pt}{9.6pt}\selectfont,firstline=227,lastline=243]{../src/black_box_06_the_verdict.py}
\end{frame}

\resultframe{The verdict, twelve times over}{../figures/black_box_06_the_verdict.png}{%
    Champions occupy a narrow \texttt{x} interval while their fitness differs by 21,504×, which is inconsistent with convergence to a finite maximum}

% ================= black_box_07_opening_the_box =================
\section{Opening the box}

\begin{frame}{Opening the box}
Six steps have treated complicated\_one() as a black box, the way the book
presents it. The searches are done; now we read the one line that was doing
all the damage. Every term of the sum is a numerator over a denominator, and
the denominator is
\end{frame}

\begin{frame}[fragile]{(1) denominator()}
\lstinputlisting[basicstyle=\ttfamily\fontsize{8pt}{9.6pt}\selectfont,firstline=75,lastline=87]{../src/black_box_07_opening_the_box.py}
\end{frame}

\begin{frame}[fragile]{(2) the pole}
\lstinputlisting[basicstyle=\ttfamily\fontsize{7pt}{8.5pt}\selectfont,firstline=98,lastline=119]{../src/black_box_07_opening_the_box.py}
\end{frame}

\begin{frame}[fragile]{(3) the ladder}
\lstinputlisting[basicstyle=\ttfamily\fontsize{7pt}{8.5pt}\selectfont,firstline=121,lastline=141]{../src/black_box_07_opening_the_box.py}
\end{frame}

\resultframe{Opening the box}{../figures/black_box_07_opening_the_box.png}{%
    The apparent optimum is a pole near \texttt{x = -10.9535836}; the reported fitness measures distance to the singularity}

\section{Next}

\begin{frame}{Next}
The apparent optimum is a pole near \(x = -10.9535836\). The reported fitness was distance to a singularity.

\vspace{0.8em}
\textbf{Lesson 09 --- Binary Combinatorial Optimisation}

The picture is gone for good. Bits encode selection, assignment, placement. Feasibility has to stay visible: penalty, repair, or an exact check.
\end{frame}
\end{document}
